\documentclass[11pt, a4paper, oneside]{article}
\usepackage[ngerman]{babel}
\usepackage{worksheet}

\begin{document}
	\author{L. Bung}
	\title{Parameter trigonometrischer Funktionen}
	\subject{Mathematik}
	\class{FHR1}
	\maketitle
	
	\task{Parameter trigonometrischer Funktionen}
	
	\begin{minipage}{.75\textwidth}
		a) Untersuchen Sie den Einfluss der Parameter $a$, $b$, $c$ und $d$ auf die Graphen der Funktionen $a \cdot \sin(b \cdot (x-c))+d$ und $a \cdot \cos(b \cdot (x-c))+d$.
		Füllen Sie dazu die Tabelle aus.
	\end{minipage}
	\begin{minipage}{.25\textwidth}
		\qrlink{https://www.geogebra.org/m/wsqrgetx}
	\end{minipage}
	
	\begin{table}[h!]
		\centering
		\begin{tabularx}{\textwidth}{|l|X|}
			\hline
			\textbf{Parameter} & \textbf{Bedeutung für den Graphen}\\
			\hline
			\hline
			$a$ & \phantom{x}\vspace{1.5cm}\\
			\hline
			$b$ & \phantom{x}\vspace{1.5cm}\\
			\hline
			$c$ & \phantom{x}\vspace{1.5cm}\\
			\hline
			$d$ & \phantom{x}\vspace{1.5cm}\\
			\hline
		\end{tabularx}
	\end{table}
	
	b) Untersuchen Sie, wie Sie die Parameter einstellen müssen, damit der Funktionsgraph von ${a \cdot \sin(b \cdot x + c) + d}$ aussieht wie der Funktionsgraph von $g(x) = \cos(x)$.
	
	Gibt es mehrere Möglichkeiten?
	
	\lines[3cm]
	
	\begin{hint}{Die Periode trigonometrischer Funktionen}
		\phantom{x}\vspace{5cm}
	\end{hint}
	
	\task{Eigenschaften trigonometrischer Funktionen}
	
	a) Geben Sie jeweils an, wie der Graph der angegebenen Funktion aus dem Graphen von $f(x) = \sin(x)$ entsteht.
	
	\begin{multicols}{3}
		\begin{enumerate}[label=(\roman*)]
			\item $a(x) = 2\sin(x)+1$
			\item $b(x) = 3\sin(2(x-2))-1$
			\item $c(x) = 2\cos(x-2)+3$
		\end{enumerate}
	\end{multicols}
	
	\checkered[5.5cm]
	
	b) Geben Sie für die Funktionen aus a) jeweils die Periode sowie den Definitions- und Wertebereich an.
	
	\checkered[3.5cm]
	
	\let\clearpage\relax
\end{document}